\documentclass[12pt,a4paper]{article}
\usepackage[utf8]{inputenc}
\usepackage[T1]{fontenc}
\usepackage[polish]{babel}
\usepackage{geometry}
\usepackage{enumitem}
\usepackage{titlesec}
\usepackage{parskip}
\usepackage{fancyhdr}
\usepackage{xcolor}
\usepackage{mdframed}
\usepackage{microtype}
\usepackage{booktabs}
\usepackage{hyperref}
\usepackage{multirow}
\usepackage{array}
\usepackage{longtable}
\usepackage{tabularx}

\geometry{margin=2.5cm}
\hypersetup{colorlinks=true, linkcolor=black, urlcolor=blue}

\pagestyle{fancy}
\fancyhf{}
\rhead{SDLC-LLM Benchmark}
\lhead{Zadanie R2 -- Refinement}
\rfoot{\thepage}

\titleformat{\section}{\large\bfseries}{}{0em}{}[\titlerule]
\titleformat{\subsection}{\normalsize\bfseries}{}{0em}{}
\titleformat{\subsubsection}{\normalsize\itshape}{}{0em}{}

\definecolor{metabg}{gray}{0.93}
\definecolor{promptbg}{RGB}{235,245,255}
\definecolor{claudecolor}{RGB}{210,105,30}
\definecolor{score3}{RGB}{220,240,220}
\definecolor{score2}{RGB}{255,243,205}
\definecolor{score1}{RGB}{255,220,200}
\definecolor{score0}{RGB}{255,200,200}
\definecolor{tpcolor}{RGB}{220,240,220}
\definecolor{fpcolor}{RGB}{255,220,200}
\definecolor{fncolor}{RGB}{255,243,205}

\newcommand{\scorebox}[1]{%
  \ifnum#1=3\colorbox{score3}{\textbf{3}}%
  \else\ifnum#1=2\colorbox{score2}{\textbf{2}}%
  \else\ifnum#1=1\colorbox{score1}{\textbf{1}}%
  \else\colorbox{score0}{\textbf{0}}\fi\fi\fi
}
\newcommand{\tp}[1]{\colorbox{tpcolor}{\small #1}}
\newcommand{\fp}[1]{\colorbox{fpcolor}{\small #1}}
\newcommand{\fn}[1]{\colorbox{fncolor}{\small #1}}

\begin{document}

% ─── STRONA TYTUŁOWA ────────────────────────────────────────────────
\begin{center}
  {\LARGE\bfseries Requirements Artifact}\\[1.5em]
  {\large \textbf{Case study:} Appointment Booking System}\\[0.3em]
  {\large \textbf{Model: Gemini 3 flash preview}  \quad }\\[0.3em]
  {\large \textbf{Tryb:} LLM-only \quad }\\[0.3em]

\end{center}

\vspace{1em}\hrule\vspace{1.5em}

% ─── PROMPT ─────────────────────────────────────────────────────────
\section{Prompt}

\begin{mdframed}[backgroundcolor=metabg, linewidth=0pt, innerleftmargin=10pt, innerrightmargin=10pt, innertopmargin=8pt, innerbottommargin=8pt]
\textbf{System prompt:}\\
\small Jesteś analitykiem wymagań. Odpowiadaj WYŁĄCZNIE poprawnym JSON. Żadnego Markdown ani komentarzy.
\end{mdframed}

\vspace{0.5em}
\begin{mdframed}[backgroundcolor=promptbg, linewidth=0pt, innerleftmargin=10pt, innerrightmargin=10pt, innertopmargin=8pt, innerbottommargin=8pt]
\textbf{User prompt (3-krokowy):}\\[0.3em]
\small
\texttt{[KROK 1]} Na podstawie opisu systemu wygeneruj \texttt{system\_goal}, role i wymagania funkcjonalne.\\
\texttt{[KROK 2]} Na podstawie opisu i wygenerowanych FR wygeneruj NFR i user stories.\\
\texttt{[KROK 3]} Na podstawie user stories wygeneruj acceptance criteria w formacie Given/When/Then.\\[0.3em]
\textbf{Opis systemu (wejście):}\\
System umożliwia użytkownikom rezerwację wizyt u specjalistów. Użytkownicy mogą przeglądać dostępne terminy, dokonywać rezerwacji oraz je anulować. Specjaliści zarządzają swoim grafikiem. System powinien zapobiegać konfliktom rezerwacji, ale w niektórych przypadkach dopuszczalne są wyjątki. Użytkownik może anulować wizytę w uzasadnionych przypadkach, jednak anulowanie może być ograniczone czasowo. Specjalista może modyfikować swój grafik, także dla już istniejących rezerwacji. System powinien obsługiwać różne role użytkowników i być skalowalny.
\end{mdframed}


\section{Cel systemu}
System służy do kompleksowego zarządzania procesem rezerwacji wizyt u specjalistów, umożliwiając użytkownikom planowanie spotkań, a specjalistom efektywne zarządzanie grafikiem pracy. Zapewnia on spójność danych poprzez automatyczne zapobieganie konfliktom terminów oraz obsługę ról o zróżnicowanych uprawnieniach.

\section{Role}
\begin{imtemize}
    \item \textbf{Użytkownik} - przeglądanie dostępnych terminów, dokonywanie rezerwacji, anulowanie własnych rezerwacji
    \item \textbf{Specjalista} - zarządzanie własnym grafikiem, modyfikacja istniejących rezerwacji, przeglądanie listy pacjentów/klientów
    \item \textbf{Administrator} - zarządzanie kontami użytkowników, konfiguracja parametrów systemu, nadzorowanie logów systemowych
\end{imtemize}


\section{Wymagania funkcjonalne}
\begin{description}[leftmargin=1.8cm, labelwidth=1.3cm, labelsep=0.5cm]
    \item[FR1] System musi umożliwiać użytkownikom przeglądanie dostępnych terminów u wybranych specjalistów w czasie rzeczywistym.
    \item[FR2] System musi pozwalać użytkownikowi na dokonanie rezerwacji wybranego wolnego terminu.
    \item[FR3] System musi umożliwiać użytkownikowi anulowanie wizyty, o ile nie został przekroczony zdefiniowany limit czasowy.
    \item[FR4] System musi automatycznie blokować możliwość rezerwacji terminu, który jest już zajęty, w celu uniknięcia konfliktów.
    \item[FR5] System musi umożliwiać specjaliście dodawanie, edycję i usuwanie terminów w swoim grafiku pracy.
    \item[FR6] System musi pozwalać specjaliście na modyfikację grafiku obejmującą terminy, na które zostały już dokonane rezerwacje.
    \item[FR7] System musi weryfikować uprawnienia użytkownika przed wykonaniem każdej operacji krytycznej (np. modyfikacja rezerwacji).
\end{description}

\section{Wymagania niefunkcjonalne}
\begin{description}[leftmargin=1.8cm, labelwidth=1.3cm, labelsep=0.5cm]
    \item[NFR1] System musi być skalowalny, aby umożliwić obsługę rosnącej liczby użytkowników i rezerwacji bez utraty wydajności.
    \item[NFR2] Czas odpowiedzi systemu podczas wyszukiwania dostępnych terminów nie może przekraczać 2 sekund przy obciążeniu do 1000 równoległych użytkowników.
    \item[NFR3] System musi zapewniać wysoką dostępność na poziomie 99,9\% w skali roku.
    \item[NFR4] Wszystkie dane osobowe i medyczne muszą być przechowywane i przesyłane w formie zaszyfrowanej, zgodnie z regulacjami RODO.
    \item[NFR5] Interfejs użytkownika musi być responsywny i poprawnie wyświetlać się na urządzeniach mobilnych oraz desktopowych.
\end{description}

\section{User stories}
\begin{description}[leftmargin=1.8cm, labelwidth=1.3cm, labelsep=0.5cm]
    \item[US1] Jako Użytkownik chcę przeglądać dostępne terminy u specjalistów, aby wybrać najbardziej dogodny dla mnie czas wizyty.
    \item[US2] Jako Użytkownik chcę zarezerwować wybrany termin, aby mieć gwarancję odbycia konsultacji.
    \item[US3] Jako Użytkownik chcę mieć możliwość anulowania rezerwacji w określonym czasie, aby zwolnić termin innym osobom, gdy nie mogę się stawić.
    \item[US4] Jako Specjalista chcę definiować i edytować swój grafik pracy, aby pacjenci widzieli tylko moje aktualne godziny dostępności.
    \item[US5] Jako Specjalista chcę mieć możliwość modyfikacji już istniejących rezerwacji, aby móc zarządzać nagłymi zmianami w moim planie dnia.
    \item[US6] Jako Administrator chcę zarządzać uprawnieniami i kontami użytkowników, aby zapewnić bezpieczeństwo i poprawność działania systemu.
\end{description}

\section{Acceptance criteria}

\begin{description}[leftmargin=1.5cm, labelwidth=1cm, labelsep=0.5cm]
\item[\textbf{AC1 -- }] 
\begin{itemize}[itemsep=2pt]
    \item \textbf{Given}: Użytkownik znajduje się na profilu wybranego specjalisty.
    \item \textbf{When}: wybiera datę z kalendarza.
    \item \textbf{Then}: system wyświetla listę wszystkich dostępnych godzin wizyt dla tego dnia.
\end{itemize}

\item[\textbf{AC2 -- }] 
\begin{itemize}[itemsep=2pt]
    \item \textbf{Given}: Użytkownik wybrał konkretny wolny termin.
    \item \textbf{When}: zatwierdza formularz rezerwacji.
    \item \textbf{Then}: system blokuje termin dla innych i wysyła potwierdzenie rezerwacji do użytkownika.
\end{itemize}

\item[\textbf{AC3 -- }] 
\begin{itemize}[itemsep=2pt]
    \item \textbf{Given}: Użytkownik posiada aktywną rezerwację, do której pozostało więcej niż 24 godziny.
    \item \textbf{When}: klika przycisk 'Anuluj rezerwację'.
    \item \textbf{Then}: system usuwa rezerwację z konta użytkownika i przywraca termin do puli dostępnych.
\end{itemize}

\item[\textbf{AC4 -- }] 
\begin{itemize}[itemsep=2pt]
    \item \textbf{Given}: Specjalista jest zalogowany w panelu zarządzania grafikiem.
    \item \textbf{When}: dodaje nowe przedziały czasowe swojej dostępności.
    \item \textbf{Then}: system aktualizuje widok kalendarza dla pacjentów zgodnie z nowymi ustawieniami.
\end{itemize}

\item[\textbf{AC5 -- }] 
\begin{itemize}[itemsep=2pt]
    \item \textbf{Given}: W systemie istnieje potwierdzona rezerwacja pacjenta.
    \item \textbf{When}: Specjalista zmienia godzinę wizyty w panelu zarządzania. 
    \item \textbf{Then}: system aktualizuje dane wizyty i wysyła automatyczne powiadomienie o zmianie do pacjenta.
\end{itemize}

\item[\textbf{AC6 -- }] 
\begin{itemize}[itemsep=2pt]
    \item \textbf{Given}: Administrator jest zalogowany do panelu administracyjnego.
    \item \textbf{When}: zmienia status konta użytkownika na 'zablokowany'.
    \item \textbf{Then}: użytkownik traci możliwość logowania się i wykonywania jakichkolwiek akcji w systemie.
\end{itemize}

\vspace{0.5cm}

\end{description}


% ─── PYTANIA UZNANE ZA ROZSTRZYGNIĘTE ──────────────────────────────
\section{Pytania uznane za rozstrzygnięte}

W tej wersji artefaktu model przyjął następujące założenia:

\begin{itemize}[leftmargin=1.5em]
  \item okno anulowania rezerwacji jest konfigurowalne przez administratora
  \item limit czasowy odpowiedzi systemu przy obiążeniu 10000 użytkowników wynosi 2 sekundy 
\end{itemize}

\newpage

% ─── METRYKI VS GOLD ────────────────────────────────────────────────
\section{Metryki porównawcze względem gold artifact}

\begin{mdframed}[backgroundcolor=metabg, linewidth=0pt, innerleftmargin=10pt, innerrightmargin=10pt, innertopmargin=8pt, innerbottommargin=8pt]
\small
\textbf{Legenda:} \tp{TP} = pokrywa element gold \quad \fp{FP} = nadmiarowe, nie ma w gold \quad \fn{FN} = brakujące względem gold\\[0.3em]
Metryki liczone na poziomie wymagań semantycznych (nie IDs), tzn.\ porównywana jest \textit{treść} wymagania z gold, nie jego numer.
\end{mdframed}

\vspace{0.8em}

\subsection{Wymagania funkcjonalne}

\begin{tabular}{p{2.5cm} p{7cm} l}
\toprule
\textbf{ID Gemini} & \textbf{Mapowanie na gold} & \textbf{Status} \\
\midrule
FR1 & Gold FR1 -- przeglądanie terminów & \tp{TP} \\
FR2 & Gold FR2 -- rezerwacja dostępnego terminu & \tp{TP} \\
FR3 & Gold FR5 -- umożliwienie anulowania wizyty & \tp{TP} \\
FR4 & Gold FR4 -- brak podwójnej rezerwacji & \tp{TP} \\
FR5 & Gold FR7 -- zarządzanie grafikiem specjalisty & \tp{TP} \\
FR6 & brak w gold -- dowolne anulowania wizyt przez specjalistów & \fp{FP} \\
FR7 & Gold FR5 -- autoryzacja użytkownika podczas modyfikacji wizyty & \tp{TP} \\
\midrule
\multicolumn{3}{l}{\textit{Brakujące z gold:}} \\
-- & Gold FR3 -- limit 3 aktywnych rezerwacji per user & \fn{FN} \\
-- & Gold FR6 -- zwalnianie terminów & \fn{FN} \\
-- & Gold FR8 -- statusy wizyt & \fn{FN} \\
-- & Gold FR9 -- dostęp do danych & \fn{FN} \\
-- & Gold FR10 -- konflikty czasowe & \fn{FN} \\
-- & Gold FR11 -- blokady dostępności & \fn{FN} \\
-- & Gold FR12 -- historia operacji & \fn{FN} \\
-- & Gold FR13 -- równoczesne operacje & \fn{FN} \\
\bottomrule
\end{tabular}

\vspace{0.5em}
\begin{tabular}{lrrrr}
\toprule
& \textbf{TP} & \textbf{FP} & \textbf{FN} & \\
\midrule
FR & 6 & 1 & 8 & \\
\midrule
\textbf{Precision} & \multicolumn{4}{l}{$6 / (6+1) = \mathbf{0{,}86}$} \\
\textbf{Recall}    & \multicolumn{4}{l}{$6 / (6+8) = \mathbf{0{,}43}$} \\
\textbf{F1}        & \multicolumn{4}{l}{$2 \times 0{,}86 \times 0{,}43 / (0{,}86+0{,}43) = \mathbf{0{,}57}$} \\
\bottomrule
\end{tabular}

\vspace{1em}
\subsection{Wymagania niefunkcjonalne}

\begin{tabular}{p{2.5cm} p{7cm} l}
\toprule
\textbf{ID Gemini} & \textbf{Mapowanie na gold} & \textbf{Status} \\
\midrule
NFR1 (skalowanie)    & Gold NFR5 -- skalowalność & \tp{TP} \\
NFR2 (czas odpowiedzi)     & Gold NFR2 -- wydajność & \tp{TP} \\
NFR3 (dostępność)        & brak bezpośredniego mapowania na gold & \fp{FP} \\
NFR4 (RODO)           & Gold NFR3 -- bezpieczeństwo  & \tp{TP} \\
NFR5 (interfejs)        & brak bezpośredniego mapowania na gold & \fp{FP} \\
\midrule
-- & Gold NFR1 -- spójność rezerwacji & \fn{FN} \\
-- & Gold NFR4 -- audyt & \fn{FN} \\
\bottomrule
\end{tabular}

\vspace{0.5em}
\begin{tabular}{lrrrr}
\toprule
& \textbf{TP} & \textbf{FP} & \textbf{FN} & \\
\midrule
NFR & 3 & 2 & 2 & \\
\midrule
\textbf{Precision} & \multicolumn{4}{l}{$3 / (3+2) = \mathbf{0{,}60}$} \\
\textbf{Recall}    & \multicolumn{4}{l}{$3 / (3+2) = \mathbf{0{,}60}$} \\
\textbf{F1}        & \multicolumn{4}{l}{$2 \times 0{,}60 \times 0{,}60 / (0{,}60+0{,}60) = \mathbf{0{,}60}$} \\
\bottomrule
\end{tabular}

\vspace{1em}
\subsection{Acceptance criteria}

\begin{tabular}{p{2.5cm} p{7cm} l}
\toprule
\textbf{ID Claude} & \textbf{Mapowanie na gold} & \textbf{Status} \\
\midrule
AC1 & brak mapowania na gold AC & \fp{FP} \\
AC2 & Gold AC1 -- poprawna rezerwacja & \tp{TP} \\
AC3 & Gold AC4 -- anulowanie wizyty $> $24h & \tp{TP} \\
AC4 & brak mapowania na gold AC & \fp{FP} \\
AC5 & brak mapowania na gold AC & \fp{FP} \\
AC6 & brak mapowania na gold AC & \fp{FP} \\
\midrule
-- & Gold AC2 -- limit 3 rezerwacji & \fn{FN} \\
-- & Gold AC3 -- brak double booking & \fn{FN} \\
-- & Gold AC5 -- anulowanie graniczne (dokładnie 24h) & \fn{FN} \\
-- & Gold AC6 -- konflikt czasowy nakładających się rezerwacji & \fn{FN} \\
\bottomrule
\end{tabular}

\vspace{0.5em}
\begin{tabular}{lrrrr}
\toprule
& \textbf{TP} & \textbf{FP} & \textbf{FN} & \\
\midrule
AC & 2 & 4 & 4 & \\
\midrule
\textbf{Precision} & \multicolumn{4}{l}{$2 / (2+4) = \mathbf{0{,}33}$} \\
\textbf{Recall}    & \multicolumn{4}{l}{$2 / (2+4) = \mathbf{0{,}33}$} \\
\textbf{F1}        & \multicolumn{4}{l}{$2 \times 0{,}33 \times 0{,}33 / (0{,}33+0{,}33) = \mathbf{0{,}33}$} \\
\bottomrule
\end{tabular}

\vspace{0.5em}
\begin{mdframed}[backgroundcolor=metabg, linewidth=0pt, innerleftmargin=10pt, innerrightmargin=10pt, innertopmargin=6pt, innerbottommargin=6pt]
\small\textbf{Komentarz do AC:} Niski ogólny wynik w seksji z AC wynikał z niedopasowania odpowiednich acceptance criteria przez model. Odpowiedź była ukierunkowana bardziej na zadania administratora oraz na zarządzanie rezerwacją wizyt. Recall 0.33 oznacza, że model pominął 4 kluczowe AC: limit rezerwacji, brak double bookingu, przypadek graniczny 24h i konflikty nakładających się terminów.
\end{mdframed}


% ─── METADANE I OCENA ───────────────────────────────────────────────
\section{Metadane i ocena}
\begin{table}[h]
% Zmieniamy tabular na tabularx, ustawiamy szerokość na \textwidth 
% i zmieniamy typ drugiej kolumny z 'l' na 'X'
\begin{tabularx}{\textwidth}{l X}
\toprule
\textbf{Pole} & \textbf{Wartość} \\
\midrule
Model & \texttt{gemini-3-flash-preview} \\
Tryb & LLM-only \\
Case study & appointment-booking (ver2 challenging) \\
Czas łącznie & 45,7s \\
\midrule
\multicolumn{2}{l}{\textit{Ocena (0--3)}} \\
\midrule
Correctness (M1)     & \scorebox{2} \quad merytorycznie poprawny, ale brak m.in. limitu 3 rezerwacji oraz statusów rezerwacji \\
Completeness (M2)    & \scorebox{1} \quad pokrywa główne założenia, ale dalej brakuje kluczowych NFR oraz FR \\
Consistency (M3)     & \scorebox{3} \quad wewnętrznie spójny, brak sprzeczności \\
Clarity (M4)         & \scorebox{3} \quad czytelna struktura, AC w formacie Given/When/Then \\
Maintainability (M5) & \scorebox{2} \quad użyteczny jako wejście dla Architecture Team po uzupełnieniu brakujących FR oraz NFR \\
\midrule
\multicolumn{2}{l}{\textit{Metryki vs gold}} \\
\midrule
FR Precision / Recall / F1   & 0.86 / 0.43 / 0.57 \\
NFR Precision / Recall / F1  & 0.60 / 0.60 / 0.60 \\
AC Precision / Recall / F1   & 0.33 / 0.33 / 0.33 \\
\bottomrule
\end{tabularx}
\end{table}


\end{document}